\subsection{The Third Week}\label{sec:week3}

The Third Sunday takes its old name from the first word of its Entrance
Antiphon, \latin{Gaudete}, ``Rejoice.'' The words are St Paul's:

\begin{quote}
Rejoice in the Lord always: again, I say, rejoice. Let your modesty be known to
all men. The Lord is nigh.\par
{\smallskip\noindent\small ---Philippians 4:4--5\par}
\end{quote}

\noindent The Missal before the reform sang them at length, with the verse that
follows them in the epistle; the present antiphon keeps a shorter form. The
Lord is near, and the Church lets the joy show. The priest may wear rose
instead of violet where that is the custom, and the calendar of the United
States gives violet or rose for this Sunday.

In Year B the first reading is the voice of the anointed one in Isaiah
(61:1--2a, 10--11): ``The spirit of the Lord is upon me, because the Lord hath
anointed me: he hath sent me to preach to the meek, to heal the contrite of
heart, and to preach a release to the captives.'' Christ read this page aloud in
the synagogue of Nazareth, folded the book, and said, ``This day is fulfilled
this scripture in your ears.'' In place of a psalm the Church sings the
Canticle of Mary, and St Paul gives the Thessalonians the rule of the day:
``Always rejoice. Pray without ceasing. In all things give thanks.'' The Gospel
is John's testimony in St John: asked who he is, he says, ``I am the voice of
one crying in the wilderness, make straight the way of the Lord, as said the
prophet Isaias.''

On the same Sunday the Office of Readings reads St Augustine on that verse. His
sermon turns on two Latin words, \latin{vox} and \latin{verbum}, voice and word:
\latin{Vox Joannes, Verbum Christus}, ``John is the voice, Christ is the Word.''
A voice carries a word from one mind to another and then dies away; the word
remains. \latin{Tolle verbum, quid est vox? Ubi nullus est intellectus, inanis
est strepitus. Vox sine verbo aurem pulsat, cor non ædificat}, ``Take away the
word, and what is the voice? Where there is no meaning, it is an empty noise. A
voice without a word strikes the ear; it does not build up the heart.'' And what
does the voice ask of us? \latin{Quid est, Viam parate; nisi, humiliter
cogitate?} ``What is `Prepare the way,' but `Think humbly'?'' And he bids his
hearers take the example of humility from John himself, the voice that was
taken for the Word and would not let himself be taken for the Christ.

In Year A Isaiah sees the desert rejoice, St James bids the Church be patient
like the farmer who ``waiteth for the precious fruit of the earth,'' and the
Lord praises John as more than a prophet. In Year C Zephaniah sings to the
daughter of Sion of the Lord in her midst: ``he will rejoice over thee with
gladness, he will be silent in his love, he will be joyful over thee in
praise.'' St Paul sounds his \latin{Gaudete}, and the crowds ask John what they
must do.

In any year John's first counsel to the crowds can be kept to the letter in the
week of Gaudete: ``He that hath two coats, let him give to him that hath none.''
A household can go through its closets together, choose the coat that is not
worn, and give it this week to someone who has none.

The weekdays of the third week run only until 16 December, and in many years
some of them are not read at all. On Monday the pagan seer Balaam, hired to
curse Israel, blesses it, and sees far off: ``I shall see him, but not now: I
shall behold him, but not near. A star shall rise out of Jacob and a sceptre
shall spring up from Israel.'' On Tuesday Zephaniah promises a remnant: ``And I
will leave in the midst of thee a poor and needy people: and they shall hope in
the name of the Lord.'' St Augustine, commenting on that passage in the
\emph{City of God}, says who they are: ``These are the remnant of that nation
who have believed in Christ.'' On Wednesday the Church hears the verse from
which Advent's chant \latin{Rorate caeli} is drawn:

\begin{quote}
Drop down dew, ye heavens, from above, and let the clouds rain the just: let the
earth be opened, and bud forth a saviour: and let justice spring up together: I
the Lord have created him.\par
{\smallskip\noindent\small ---Isaiah 45:8\par}
\end{quote}

\noindent In the Gospel the disciples of John come to ask, ``Art thou he that art
to come?'' and the Lord points to the blind who see and the poor who have the
Gospel preached to them. In the Office that day St Irenaeus says that God's Word
``in the last times was made a man among men, that He might join the end to the
beginning, that is, man to God.''

The chant \latin{Rorate caeli desuper} takes Wednesday's verse for its
refrain, and its other verses are drawn from Isaiah's lament and comfort. It is
given in Latin and English in Appendix~\ref{sec:prayers}, and a household can
sing or say it on the evenings of this week.

St Lucy, virgin and martyr of Syracuse in the persecution of Diocletian, is
kept on 13 December. The Roman Martyrology of 1902 told how the men who meant to
drag her away to be dishonoured could not stir her:
\latin{nullatenus eam movere potuerunt, nec funibus additis, nec boum jugis
plurimis}, ``they could in no way move her, neither by adding ropes nor by many
yoke of oxen.'' She died at last by the sword.

On 14 December comes St John of the Cross, Carmelite, poet and Doctor of the
Church. Among his ballads are \emph{The Desires of the Holy Fathers},
\emph{The Incarnation} and \emph{The Nativity}, and the first gives the prayers
of the holy men of old who waited for Christ. It sounds like the Advent
Lectionary set to verse:

\begin{verse}
`Hasten, then, Thy work, and finish;\\
Send Him, Lord, Whom Thou wilt send,'\\
Was the cry of one. Another's,\\
`O that He the heavens would rend!'\\
`That I might behold His coming,\\
And my wail be turned to mirth;\\
Let the clouds rain down the Just One,\\
So long desired on the earth;'
\end{verse}

\noindent The prayer of Isaiah on the First Sunday and the \latin{Rorate} of this
week are both there. And among his short sayings is the one that measures every
Advent: ``At the close of life you will be examined as to your love: learn then
to love God as He wishes to be loved, and give up all that is your own.''
